\documentclass[10pt,a4paper]{amsart}
\usepackage[margin=19mm]{geometry}
\usepackage{amsmath,amssymb,mathtools}
\usepackage[scr=rsfs,cal=euler]{mathalfa}
\usepackage{xfrac}
\usepackage[unicode,hidelinks,bookmarks=false]{hyperref}
\newcommand{\ie}{i.e.\ }
\newcommand{\cf}{cf.\ }
\newcommand{\resp}{respectively\ }
\newcommand{\Subsection}{\subsection}
\providecommand{\nobreakdash}{\nobreak\hbox{-}}
\setlength{\emergencystretch}{2em}
\multlinegap0pt
\usepackage{csquotes}
\usepackage{mathtools}
\usepackage{bm}
\newtheorem{lemma}{Lemma}
\usepackage{cleveref}
\crefname{lemma}{Lemma}{Lemmas}
\title{Product Structure and Treewidth of Hyperbolic Uniform Disk Graphs}
\author{Thomas Bl\"asius, Emil Dohse, Deborah Haun, Laura Merker}
\date{Source: arXiv:2603.18997v1 (CC BY 4.0)}
\begin{document}
\maketitle

\noindent\emph{Source and licence.} Product Structure and Treewidth of Hyperbolic Uniform Disk Graphs. Version arXiv:2603.18997v1 (CC BY 4.0), distributed by arXiv under the Creative Commons Attribution 4.0 International licence, http://creativecommons.org/licenses/by/4.0/, as recorded in the arXiv licence metadata for this exact version and shown on the abstract page https://arxiv.org/abs/2603.18997v1. Full licence notice: \texttt{licenses/arxiv-2603.18997-CC-BY-4.0.txt}.\par\smallskip

\noindent\textit{Editorial scope.} Counted unit: the article's lemma lem:Gr\_clique-levels (source0.tex lines 1041--1044). The article's complete original proof of it is delivered with it (source0.tex lines 1046--1083, one \texttt{proof} environment of 38 lines). No mathematics is written, repaired or re-derived: the statement and the proof are the article's own lines, in the article's own notation, and no retained line is substituted or re-flowed.  No further source-local context is needed: every cross-reference of every retained line resolves inside the two delivered spans.  Nothing was omitted from any retained span, so the package carries no omission record.  No retained line contains a citation, so the wrapper carries no bibliography and no bibliography entry.  No retained line contains a figure environment, an image-inclusion command or drawing code, so no image is delivered and the package declares no asset dependency.  Editorial formatting only, in addition: an AMS \texttt{amsart} preamble that restates the article's own declarations and macros used by the retained lines, namely the environments \texttt{lemma} on separate counters, together with the standard packages those lines need.  The retained environments print in this document as Lemma 1 in order of appearance, whereas the article prints the same items with the numbering of its own full document.  The complete native source of the article as distributed in the arXiv e-print for this exact version is bundled unchanged as \texttt{sources/196746/source0.tex}.
\begin{lemma}\label{lem:Gr_clique-levels}
    Let $ r \geq 1 $, $ c \geq 0 $ be integers, and let $ S \subseteq [r] $ be a set such that for every three elements $ i < j < k $ of $ S $, we have $ k - j \leq j - i + c $.
    Then $ |S| < \log_2(r) + c + 2 $.
\end{lemma}

\begin{proof}%[Proof of \cref{lem:Gr_clique-levels} (greedy)]
    Let us first describe how an optimal set $ S^* $ looks like before arguing that this is indeed optimal.
    We choose the elements for $ S^* $ greedily, adding always the largest element such that $ k - j \leq j - i + c $ is maintained for any three elements $ i < j < k $ in~$ S^* $.
    We claim that $ S^* $ contains roughly the $ c $ largest elements and starting from there, leaves exponentially growing gaps when smaller elements are added.
    More formally, we define $ S^* $ to contain the elements $ r, \dots, r - c $ and the elements $ s_a^* = r - c - 2^a $ for $ a = 0, \dots, \lfloor \log_2(r - c - 1) \rfloor $, and then show that $ S^* $ indeed results from choosing the elements greedily.
    % $ s_a^* = s_{a-1}^* - 2^a $ for $ a = 0, \dots, \log_2(r) $, where $ s_{-1}^* = r - (c+1) $.
    To show that $ S^* $ is a valid choice for $ S $, it suffices to compare each two consecutive elements with the largest element $ r $, where consecutive refers to two elements $ s, s' \in S^* $ such that there is no third element $ x \in S^* $ with $ s < x < s' $.
    Indeed, the largest $ c + 2 $ elements may be contained in $ S^* $ as the three elements $ r, r-c, r - c - 2^0 $ satisfy $ r - (r - c) = c \leq c+1 = (r-c) - (r-c - 2^0) + c $.
    Similarly, for $ a > 0 $, we have 
    $ r - s_{a-1}^* 
    = r - (r - c - 2^{a-1}) 
    \leq 2^{a-1} + c
    \leq 2^{a} - 2^{a-1} + c
    \leq (r - c - 2^{a-1}) - (r - c - 2^{a}) + c
    = s_{a-1}^* - s_{a}^* + c$.
    Since the largest $ c + 2 $ elements form an interval and for the $ s_a^* $, all inequalities are tight, none of the elements can be chosen larger without violating the requirements.
    Thus, $ S^* $ indeed is the result of a greedy approach.
    Moreover, the resulting set has size $ |S^*| = c + 1 + \lfloor \log_2(r - c - 1) \rfloor + 1 < \log_2(r) + c + 2 $.

    It remains to show that there is no larger set $ S $ meeting the requirements. 
    To do so, consider all subsets of $ [r] $ with maximum cardinality satisfying $ k - j \leq j - i + c $ and let $ S $ be the one that maximizes the sum of its elements.
    We show that $ S = S^* $, which means that $ S^* $ indeed has maximum cardinality.
    Assume for contradiction that $S \neq S^*$ and let $ t \in [r] $ be the largest integer such that $ S $ and $ S^* $ disagree about $ t $. 
    Since $ S^* $ is constructed greedily from large to small elements and the two sets agree on all larger elements, we have that if $ t \in S $, then also $ t \in S^* $.
    Hence, $ t \in S^* $ but $ t \notin S $.
    
    We finish the proof by moving some elements of $ S $ so that the sum gets larger, which contradicts the choice of $ S $.
    First, if $ t = r $, then add 1 to all elements of $ S $.
    This does not change the differences, and thus yields a valid set of the same cardinality whose sum is larger.
    So assume $ t < r $ and let $ t' \in S $ be the largest integer with $ t' < t $.
    Note that $ t' $ exists as $ S $ is of maximum cardinality and thus at least as big as $ S^* $.
    Now, we claim that the set $ S' $ obtained from $ S $ by replacing $ t' $ with  $t$ satisfies $ k - j \leq j - i + c $ for every $ i < j < k $.
    We only need to check triples containing $ t $.
    Also recall that we may assume $ k = r $ as this maximizes $ k - j $.
    First, for $ j = t $, the difference between $ i $ and $ j $ is now even larger, while the difference between $ j $ and $ k $ shrinks, so $ k - j = k - t < k - t' \leq t' - i < t - i = j - i $ for every $ i < t' < t$ in $ S $, and thus also in $ S' $.
    Second, if $ i = t $, then the $ i, j, k $ are all contained in $ S^* $ by the choice of $ t $ and thus satisfy the requirement for $ S' $.
    We conclude that $ S = S^* $ and therefore $ |S| = |S^*| < \log_2(r) + c + 2 $.
\end{proof}

\end{document}
